\section{Model and honest interval criterion}

One observed record is \(O=(X,A,Y)\) in
\[
\Omega=[0,1]\times\{0,1\}\times\{0,1\},
\]
equipped with its product Borel structure. The data
\(\mathcal O=(O_1,\ldots,O_n)\) are independent draws from an observed law
\(P\); sample sizes are positive integers, and \(A,Y\in\{0,1\}\). Write
\(p_{ay}(P,x)=P(A=a,Y=y\mid X=x)\). The observed-law carrier used below
consists of normalized conditional four-cell tables whose entries are
continuous and strictly positive on \([0,1]\), with a covariate marginal of
full support. Thus all conditional quantities have pointwise continuous
versions throughout the covariate interval.

Define the arm risks and treatment probability by
\[
\mu_a(P,x)=\frac{p_{a1}(P,x)}{p_{a0}(P,x)+p_{a1}(P,x)},
\qquad
e(P,x)=p_{10}(P,x)+p_{11}(P,x),
\]
and define the native logits
\[
g(P,x)=\operatorname{logit}\{e(P,x)\},
\qquad
\nu(P,x)=\operatorname{logit}\{\mu_0(P,x)\}.
\]
For \(0<\gamma\le1\), our Hölder convention is
\[
[f]_\gamma=\sup_{x\ne z}\frac{|f(x)-f(z)|}{|x-z|^\gamma}.
\]
The public exponent domain is
\(\mathcal D=\{(\alpha,\beta):0<\beta<1/4,\ \beta<\alpha\le1\}\).
Generic positive constants may change from line to line; \(\|\cdot\|_\infty\)
denotes the supremum norm, \(\lesssim\) hides a positive constant independent
of sample size and evaluation radius, and \(\asymp\) denotes two such bounds.

The comparison output space \(\mathcal C\) consists of every nonempty connected
Borel subset of \([-1/2,1/2]\). Encode an element by its lower and upper
endpoints together with two endpoint-inclusion indicators, excluding the empty
encoding; equip this space with the induced product Borel structure. The
attaining construction below returns the closed member \([l,u]\) with rational
endpoints. Randomized procedures also receive an independent seed
\(U\sim\operatorname{Unif}[0,1]\).
This specifies the record, sampling, and output spaces before the model and
honesty conditions.

The design law is known and fixed as follows.

\begin{assumptionv}{ass:uniform-design}[Uniform covariate design]
The covariate \(X\) has the uniform probability law \(\lambda\) on the unit interval under the observed law \(P\):
\[
P\circ X^{-1}=\lambda.
\]
\end{assumptionv}

The known uniform design is a fixed-density specialization of the known-density setting \citep{McCleanBalakrishnanKennedyWasserman2024DCDR}. It fixes the covariate integration measure across observed laws, allowing the model restrictions to focus on treatment assignment and conditional outcome risks.

We express those risks on the log-odds scale. The log-odds map \leanref{sym:\operatorname{logit}}{\ensuremath{\operatorname{logit}}} is defined first, followed by the native control-risk logit.

\begin{definitionv}{synth_2}[Log-odds map definition]
The log-odds map \(\operatorname{logit}\colon\mathbb R\to\mathbb R\) is defined by
\[
\operatorname{logit}(t)=\log\!\left(\frac{t}{1-t}\right),
\qquad t\in\mathbb R.
\]
Here division by zero has value zero, and the real logarithm is extended by \(\log 0=0\) and \(\log t=\log|t|\) for \(t<0\). For \(0<t<1\), this is the usual log-odds transformation.
\end{definitionv}

The conditional risks in \cref{obj:def:model} lie in \((0,1)\), where this transformation has its usual statistical meaning. The native prognosis logit \leanref{sym:\nu}{\ensuremath{\nu(P,x)}} applies it to the control-arm risk.

\begin{definitionv}{synth_4}[Native prognosis logit]
The native prognosis logit \(\nu(P,x)\) is the log odds of the conditional outcome risk under control, for \(x\in[0,1]\) and an observed probability law \(P\) with binary treatment and outcome whose conditional four-cell table satisfies:
\begin{itemize}
\item \textbf{(Conditional probabilities.)} At each covariate value, the four cells sum to one and give a conditional representation of \(P\) given the covariate.
\item \textbf{(Continuity and interiority.)} Each cell probability is continuous in the covariate and lies strictly between zero and one at every covariate value.
\item \textbf{(Full support.)} The covariate marginal assigns positive probability to every nonempty relatively open subset of \([0,1]\).
\end{itemize}
Its defining formula is
\[
\nu(P,x)
:=
\operatorname{logit}\!\left(
\frac{\Pr_P(\text{treatment}=0,\text{outcome}=1\mid X=x)}
{\Pr_P(\text{treatment}=0,\text{outcome}=0\mid X=x)
+\Pr_P(\text{treatment}=0,\text{outcome}=1\mid X=x)}
\right),
\qquad
\operatorname{logit}(t):=\log\!\left(\frac{t}{1-t}\right).
\]
\end{definitionv}

Consequently, \(\nu(P,x)=\operatorname{logit}\{\mu_0(P,x)\}\). The logistic map \leanref{sym:\ell}{\ensuremath{\ell}} converts a real argument \leanref{sym:t}{\ensuremath{t}} back to a probability.

\begin{definitionv}{synth_1}[The logistic map]
The logistic map \(\ell:\mathbb R\to\mathbb R\) is defined by
\[
\ell(t)=\frac{1}{1+\exp(-t)},\qquad t\in\mathbb R.
\]
\end{definitionv}

On interior probabilities, the logistic and log-odds maps are inverses. We can therefore describe treatment assignment through its native logit and fix the scalar outcome contrast before imposing homogeneity.



The public smoothness exponents govern regularity of these native logits. We fix their notation and the seminorm convention before stating the bounds.



The model combines a homogeneous conditional effect with fixed scalar and nuisance-function envelopes. Its structural restriction is the following logistic relation.

\begin{assumptionv}{ass:homogeneous-logit}[Homogeneous conditional log-odds relation]
The treated-arm conditional outcome risk \(\mu_1(P,x)\) satisfies
\[
\mu_1(P,x)=\ell\{\nu(P,x)+\theta(P)\}
\qquad (x\in[0,1]),
\]
where \(\ell\) is the logistic map, \(\nu(P,x)\) is the native prognosis logit, and \(\theta(P)\) is the homogeneous conditional log-odds coefficient.
\end{assumptionv}

The logistic partially linear coefficient with binary exposure is a standard specification \citep{Tan2019LogisticDR}. Here it makes the conditional log-odds contrast constant across covariate values:
\[
\operatorname{logit}\{\mu_1(P,x)\}
-
\operatorname{logit}\{\mu_0(P,x)\}
=\theta(P).
\]
Equivalently, the treated-to-control conditional odds ratio is \(\exp\{\theta(P)\}\). Thus the coefficient describes a common conditional odds contrast while allowing the baseline risk to vary with the covariate.

Under consistency and conditional independence of treatment and potential outcomes given the covariate, this contrast has the corresponding causal interpretation \citep{Rubin1974,Rosenbaum1983}. The \leanref{S-4}{full-data representation} associated with the observed law is specified in \cref{obj:def:causal-extension}; \cref{obj:lem:observable-odds} supplies the identification result and its causal interpretation in the appendix.

We maintain a compact range for the common coefficient.

\begin{assumptionv}{ass:effect-envelope}[Compact scalar effect envelope]
For an observed law \(P\), the log-odds contrast at covariate zero,
\[
\theta(P)=\operatorname{logit}\bigl(\mu_1(P,0)\bigr)-\nu(P,0),
\]
satisfies
\[
|\theta(P)|\le \frac12.
\]
\end{assumptionv}

A compact coefficient region is standard in logistic coefficient analysis \citep{LiuZhangZhou2020LogisticDML}. The numerical bound \(1/2\) is maintained here as part of the model and gives a common effect range for the interval comparison.

The treatment mechanism is controlled by a separate logit envelope.

\begin{assumptionv}{ass:propensity-envelope}[Propensity logit envelope]
The native propensity-logit function \(g(P,\cdot)\) satisfies
\[
\|g(P,\cdot)\|_\infty\le 1.
\]
\end{assumptionv}

The bounded propensity logit is a quantitative bounded-overlap specialization, with the constant \(1\) maintained here \citep{RobinsTchetgenLiVanderVaart2009Minimax}. Since \(e(P,x)=\ell\{g(P,x)\}\), it places treatment probabilities between \(\ell(-1)\) and \(\ell(1)\) uniformly over covariate values.

A corresponding envelope controls the baseline outcome risk.

\begin{assumptionv}{ass:prognosis-envelope}[Prognosis logit envelope]
The native prognosis-logit function \(\nu(P,\cdot)\) satisfies
\[
\|\nu(P,\cdot)\|_\infty\le 1.
\]
\end{assumptionv}

The bounded prognostic logit belongs to the bounded logistic framework, with this numerical envelope maintained for the present analysis \citep{LiuZhangZhou2020LogisticDML}. Together with \cref{obj:ass:effect-envelope}, it keeps both arm-specific outcome risks uniformly in the interior of the probability interval.

The remaining nuisance restrictions describe variation across covariate values, beginning with treatment assignment.

\begin{assumptionv}{ass:propensity-holder}[Propensity Hölder seminorm bound]
The Hölder seminorm of the native propensity-logit function \(g(P,\cdot)\) at the public smoothness exponent \(\alpha\) satisfies
\[
[g(P,\cdot)]_\alpha\le 2.
\]
\end{assumptionv}

A Hölder ball provides the standard smoothness restriction on the treatment mechanism \citep{RobinsTchetgenLiVanderVaart2009Minimax}. Here the ball is imposed on the native propensity logit, with prescribed radius \(2\); its exponent controls how rapidly treatment assignment can vary over nearby covariate values.

Baseline prognosis has its own public exponent.

\begin{assumptionv}{ass:prognosis-holder}[Prognosis Hölder bound]
The native prognosis-logit function \(\nu(P,\cdot)\) satisfies the Hölder seminorm bound
\[
[\nu(P,\cdot)]_\beta\le 2.
\]
\end{assumptionv}

The prognostic Hölder ball is likewise a standard regularity condition, imposed here on the control-risk logit with prescribed radius \(2\) \citep{LiuMukherjeeRobinsTchetgen2021Adaptive}. Distinct exponents permit treatment assignment and baseline prognosis to have different degrees of smoothness.

We collect the structural relation and the quantitative restrictions into two classes. The homogeneous logistic class \leanref{sym:\mathcal L}{\ensuremath{\mathcal L}} records the common conditional effect relation.

\begin{definitionv}{def:logistic-class}[Homogeneous logistic class \(\mathcal L\)]
\[
\mathcal L
=
\left\{
P:\mu_1(P,x)=\ell\{\nu(P,x)+\theta(P)\}
\ \text{for every }x\in[0,1]
\right\}.
\]
Here \(\mu_1(P,x)\) is the treated-arm conditional outcome risk, \(\ell\) is the logistic map, \(\nu(P,x)\) is the native prognosis logit, and \(\theta(P)\) is the homogeneous conditional log-odds coefficient. The defining relation is that of \cref{obj:ass:homogeneous-logit}.
\end{definitionv}

For a fixed public exponent pair, the ambient model \leanref{sym:\mathcal M}{\ensuremath{\mathcal M(\alpha,\beta)}} adds the known design law and the prescribed envelopes and seminorm bounds.

\begin{definitionv}{def:model}[Ambient model \(\mathcal M(\alpha,\beta)\)]
\(\mathcal M(\alpha,\beta)\) is the class of observed laws \(P\) satisfying the following conditions, for each \((\alpha,\beta)\in\mathcal D\), the prescribed domain of public smoothness exponents:
\begin{itemize}
\item \textbf{(Uniform design.)} With \(\lambda\) denoting the uniform probability law on the unit interval,
\[
P\circ X^{-1}=\lambda.
\]
See \cref{obj:ass:uniform-design}.
\item \textbf{(Homogeneous logit.)} The treated-arm conditional outcome risk \(\mu_1(P,x)\), logistic map \(\ell\), native prognosis logit \(\nu(P,x)\), and homogeneous conditional log-odds coefficient \(\theta(P)\) satisfy
\[
\mu_1(P,x)=\ell\{\nu(P,x)+\theta(P)\}
\qquad (x\in[0,1]).
\]
See \cref{obj:ass:homogeneous-logit}.
\item \textbf{(Effect envelope.)}
\[
|\theta(P)|\le \frac12.
\]
See \cref{obj:ass:effect-envelope}.
\item \textbf{(Propensity envelope.)} The native propensity-logit function \(g(P,\cdot)\) satisfies
\[
\|g(P,\cdot)\|_\infty\le 1.
\]
See \cref{obj:ass:propensity-envelope}.
\item \textbf{(Propensity smoothness.)}
\[
[g(P,\cdot)]_\alpha\le 2.
\]
See \cref{obj:ass:propensity-holder}.
\item \textbf{(Prognosis envelope.)}
\[
\|\nu(P,\cdot)\|_\infty\le 1.
\]
See \cref{obj:ass:prognosis-envelope}.
\item \textbf{(Prognosis smoothness.)}
\[
[\nu(P,\cdot)]_\beta\le 2.
\]
See \cref{obj:ass:prognosis-holder}.
\end{itemize}
The bracket notation denotes the Hölder seminorm at the indicated public exponent.
\end{definitionv}

This ambient class determines the coverage requirement. To evaluate precision near different effect magnitudes, we introduce an evaluation radius \leanref{sym:r}{\ensuremath{r}} in \([0,1/2]\).

\begin{assumptionv}{ass:evaluation-radius}[Effect evaluation radius]
The homogeneous conditional log-odds coefficient \(\theta(P)\) satisfies
\[
|\theta(P)|\le r.
\]
\end{assumptionv}

The nested evaluation subregion follows the ambient-coverage and subclass-length criterion of the honest-interval literature \citep{CaiLow2004ConfidenceAdaptation}. The radius indexes the laws over which expected length is assessed, while the ambient class remains the coverage domain.

The effect-radius slice \leanref{sym:\mathcal M_r}{\ensuremath{\mathcal M_r(\alpha,\beta)}} makes this evaluation class explicit.

\begin{definitionv}{def:radius-model}[Effect-radius slice \(\mathcal M_r(\alpha,\beta)\)]
\[
\mathcal M_r(\alpha,\beta)
=
\{P\in\mathcal M(\alpha,\beta):|\theta(P)|\le r\}
\qquad
\bigl((\alpha,\beta)\in\mathcal D,\ r\in[0,1/2]\bigr).
\]
Here \(\mathcal D\) is the prescribed domain of public smoothness exponents and \(\theta(P)\) is the homogeneous conditional log-odds coefficient. This defines the slice of the ambient model in \cref{obj:def:model} satisfying \cref{obj:ass:evaluation-radius}.
\end{definitionv}

These slices are nested in the radius. At \(r=0\), they contain the ambient laws with zero conditional log-odds contrast; at \(r=1/2\), the slice equals the ambient model. The nuisance restrictions remain fixed throughout this comparison.

Before stating honesty, we specify the sample, interval outputs, and independent randomization used by the comparison class.



Including an independent seed permits randomized procedures in the comparison class. A deterministic procedure is represented by a map that is constant in its seed argument. Coverage probabilities and expected lengths average over both the original records and this independent randomization.

We require coverage uniformly over the ambient model.

\begin{assumptionv}{ass:global-coverage}[Global interval coverage]
The interval procedure \(I(\mathcal O,U)\), evaluated on the tuple of observed records \(\mathcal O\) and an independent uniform randomization seed \(U\), satisfies
\[
\inf_{P\in\mathcal M(\alpha,\beta)}
(P^{\otimes n}\otimes\lambda)
\{\theta(P)\in I(\mathcal O,U)\}\ge 0.90,
\]
where \(\mathcal M(\alpha,\beta)\) is the ambient observed-law model, \(\lambda\) is the uniform probability law on the unit interval, and \(\theta(P)\) is the homogeneous conditional log-odds coefficient.
\end{assumptionv}

Global honesty with subclass expected-length evaluation is the standard comparison criterion, specialized here to the binary treatment and outcome experiment \citep{CaiLow2004ConfidenceAdaptation}. It is a procedure-admissibility requirement: for each public exponent pair and sample size, an admissible interval covers the coefficient with probability at least \(90\%\) at every law in the ambient model.

The honest procedure class \leanref{sym:\mathcal A}{\ensuremath{\mathcal A_n(\alpha,\beta)}} collects the maps satisfying this requirement.

\begin{definitionv}{def:honest-procedures}[Honest interval procedures $\mathcal A_n(\alpha,\beta)$]
For every \(n\ge1\) and \((\alpha,\beta)\in\mathcal D\), define
\[
\mathcal A_n(\alpha,\beta)
=
\left\{
I:\Omega^n\times[0,1]\to\mathcal C:
\inf_{P\in\mathcal M(\alpha,\beta)}
(P^{\otimes n}\otimes\lambda)
\{\theta(P)\in I(\mathcal O,U)\}
\ge0.90
\right\},
\]
where the maps \(I\) are Borel and have measurable coverage events and interval length, as required by \cref{obj:ass:global-coverage}. Here \(\Omega\) is the measurable space of one observed record, \(\mathcal O\) is the tuple of original observed records, \(\mathcal C\) is the space of admissible nonempty connected Borel effect intervals, and \(\lambda\) is the uniform probability law on \([0,1]\). The variable \(U\) is an independent randomization seed, \(\mathcal D\) is the prescribed domain of public smoothness exponents, and \(\theta(P)\) is the homogeneous conditional log-odds coefficient.
\end{definitionv}

For a fixed sample size and public exponent pair, this same admissible class is used at every evaluation radius. Precision is measured by the minimax expected length \leanref{sym:J}{\ensuremath{J_n(\alpha,\beta;r)}} over the corresponding slice.

\begin{definitionv}{def:length-objective}[Minimax expected length $J_n(\alpha,\beta;r)$]
The minimax expected interval length on the effect-radius slice, subject to honesty over the ambient model, is
\[
J_n(\alpha,\beta;r)
=
\inf_{I\in\mathcal A_n(\alpha,\beta)}
\sup_{P\in\mathcal M_r(\alpha,\beta)}
E_{P^{\otimes n}\otimes\lambda}
\bigl|I(\mathcal O,U)\bigr|.
\]
Here \(\mathcal A_n(\alpha,\beta)\) is the class of Borel interval procedures with global \(90\%\) coverage, \(\mathcal M_r(\alpha,\beta)\) is the effect-radius slice, \(\mathcal O\) is the tuple of original observed records, and \(U\) is an independent seed with uniform law \(\lambda\) on \([0,1]\). The quantity \(\lvert I(\mathcal O,U)\rvert\) is the length of the returned interval.
\end{definitionv}

The objective combines ambient coverage with local evaluation of expected length. In particular, evaluation at the null retains coverage over all effects in the ambient envelope. Varying the radius therefore measures how much precision an ambiently honest procedure can attain across nested effect regions.
